\documentclass[../../../uem_mat_mei_q.tex]{subfiles}

\begin{document}
    平面上に3点A，B，Cを頂点とする三角形ABCがあり，辺ABの長さは$3\sqrt{3}$，%
    辺ACの長さは$5\sqrt{3}$，$\angle\mathrm{BAC}$は$120^\circ$である．%
    三角形ABCの内心をIとし，三角形ABCの外心をOとする．

    以下の問に答えなさい．空欄内の各文字に当てはまる数字を所定の解答欄にマークしなさい．%
    ただし，分数はすべて既約分数にしなさい．%
    根号を伴う空欄は，根号の中に現れる自然数が最小となる形で答えなさい．

    \begin{enumerate}
        \item%
            辺BCの長さは$\myboxb{タ}\sqrt{\,\myboxb{チ}\,}$であり，三角形ABCの外接円%
            の直径は$\myboxb{ツテ}$である．
        \item%
            三角形ABCの面積は$\myfracc{\myboxb{トナ}\sqrt{\,\myboxb{ニ}\,}}{\myboxb{ヌ}}$であり，%
            三角形ABCの内接円の半径は$\myfracc{\myboxb{ネ}}{\myboxb{ノ}}$である．
        \item%
            $\overrightarrow{\mathrm{AI}}=%
            \myfracc{\myboxb{ハ}}{\myboxb{ヒ}}\overrightarrow{\mathrm{AB}}+%
            \myfracc{\myboxb{フ}}{\myboxb{ヘ}}\overrightarrow{\mathrm{AC}}$ %
            である．
        \item%
            $\overrightarrow{\mathrm{AO}}=%
            \myfracc{\myboxb{ホマ}}{\myboxb{ミ}}\overrightarrow{\mathrm{AB}}+%
            \myfracc{\myboxb{ムメ}}{\myboxb{モヤ}}\overrightarrow{\mathrm{AC}}$ %
            である．
    \end{enumerate}
\end{document}
